\section{Resolusi Query dan Rantai Penalaran Deduktif}

Tujuan utama membangun basis pengetahuan adalah untuk dapat menjawab pertanyaan (\textit{query}) secara otomatis melalui penalaran deduktif. Proses ini menggabungkan fakta-fakta dari $\Delta$ dan aturan-aturan dari $\Gamma$ menggunakan aturan inferensi untuk menurunkan kesimpulan baru \cite{rosen2019}.

\subsection{Proses Resolusi Query}

Resolusi query mengikuti langkah-langkah sistematis:

\begin{enumerate}
    \item \textbf{Terima query}: Identifikasi apa yang ingin dibuktikan, misalnya $Kakek(Hasan, Ali)$?
    \item \textbf{Cari fakta langsung}: Apakah query sudah merupakan fakta dalam $\Delta$? Jika ya, jawab ``benar''.
    \item \textbf{Cari aturan yang relevan}: Identifikasi aturan dalam $\Gamma$ yang konsekuennya cocok dengan query.
    \item \textbf{Instansiasi aturan}: Terapkan UI untuk mengikat variabel dengan konstanta yang sesuai.
    \item \textbf{Evaluasi anteseden}: Periksa apakah anteseden aturan terpenuhi (mungkin memerlukan penalaran rekursif).
    \item \textbf{Terapkan Modus Ponens}: Jika anteseden bernilai benar, simpulkan konsekuen.
\end{enumerate}

\subsection{Rantai Penalaran (Chaining)}

Dalam banyak kasus, pembuktian suatu query memerlukan rangkaian langkah inferensi yang saling terhubung. Terdapat dua strategi utama:

\begin{itemize}
    \item \textbf{Forward Chaining} (penalaran maju): Mulai dari fakta yang diketahui, terapkan aturan untuk menurunkan fakta-fakta baru hingga query terjawab. Strategi ini bersifat \textit{data-driven}.
    \item \textbf{Backward Chaining} (penalaran mundur): Mulai dari query yang ingin dijawab, cari aturan yang konsekuennya cocok, lalu coba buktikan anteseden-nya secara rekursif. Strategi ini bersifat \textit{goal-driven}.
\end{itemize}

\subsection{Contoh Rantai Penalaran}

\begin{contoh}
\textbf{Basis Pengetahuan:}

Fakta ($\Delta$):
\begin{align*}
&Manusia(Marcus), \quad Penduduk(Marcus, Pompeii)
\end{align*}

Aturan ($\Gamma$):
\begin{align*}
&R_1: \forall x\, (Manusia(x) \rightarrow Mortal(x)) \\
&R_2: \forall x\, (Penduduk(x, Pompeii) \rightarrow Tewas(x, Erupsi, 79)) \\
&R_3: \forall x\, (Mortal(x) \wedge Tewas(x, Erupsi, 79) \rightarrow Legenda(x))
\end{align*}

\textbf{Query}: $Legenda(Marcus)$?

\textbf{Penalaran (Forward Chaining):}
\begin{enumerate}
    \item Dari $Manusia(Marcus)$ dan $R_1$ (UI, $x = Marcus$): $Mortal(Marcus)$.
    \item Dari $Penduduk(Marcus, Pompeii)$ dan $R_2$ (UI, $x = Marcus$): $Tewas(Marcus, Erupsi, 79)$.
    \item Dari $Mortal(Marcus) \wedge Tewas(Marcus, Erupsi, 79)$ dan $R_3$ (UI, $x = Marcus$): $Legenda(Marcus)$.
    \item Query terjawab: $Legenda(Marcus)$ bernilai \textbf{Benar}. $\blacksquare$
\end{enumerate}
\end{contoh}

\begin{konsep}
Perhatikan bahwa rantai penalaran di atas menunjukkan bagaimana fakta-fakta baru ($Mortal(Marcus)$, $Tewas(Marcus, Erupsi, 79)$, $Legenda(Marcus)$) tidak secara eksplisit ada dalam KB, tetapi dapat \textbf{diturunkan} melalui inferensi logis. Inilah kekuatan utama pendekatan basis pengetahuan: kemampuan menemukan implikasi tersembunyi dari informasi yang tersedia.
\end{konsep}
